\documentclass[11pt]{article}
\usepackage[T1]{fontenc}
\usepackage{textcomp}
\usepackage[margin=0.75in]{geometry}
\usepackage{amsmath,amssymb,amsthm}
\usepackage{array,booktabs}
\usepackage{hyperref}
\setlength{\parindent}{0pt}
\newtheorem{theorem}{Theorem}
\theoremstyle{definition}
\newtheorem{definition}{Definition}
\title{Flow Proof Artifact: prop-14-gnomon-complements-equal}
\author{Generated by \texttt{flow doc proof}}
\date{Flow Proof Book}
\begin{document}
\maketitle
The complements of the parallelograms about the diameter of a gnomon are equal to one another.\\[0.5em]
\textbf{Source.} euclid — Elements, Book II, Proposition 14\\[0.5em]
\bigskip
\noindent{\large \textbf{Derived fact 1.} \textit{The complements of the parallelograms about the diameter of a gnomon are equal to one another} \hfill the Euclidean plane \cdot Euclid Book II \cdot Proposition 14: the complements about the diameter of a gnomon are equal}\par
\smallskip\noindent\textit{Coordinate: the Euclidean plane \cdot Euclid Book II \cdot Proposition 14: the complements about the diameter of a gnomon are equal}\par
\smallskip\noindent\textit{Source: euclid — Elements, Book II, Proposition 14}\par
\smallskip\noindent\textit{Built on: proposition 34: complements of parallelograms about the diameter are equal, for Book I of the Elements in on the Euclidean plane, proposition 11: a gnomon equals the rectangle by the segments, for Euclid Book II on the Euclidean plane}\par
\medskip
\noindent\textbf{Goal.} The complements of the parallelograms about the diameter of a gnomon are equal to one another\par
\noindent$\displaystyle \text{complement 1} = \text{complement 2}$\par
\medskip
\noindent
\begin{tabular}{@{} >{\raggedright\arraybackslash}p{0.06\textwidth} >{\raggedright\arraybackslash}p{0.47\textwidth} >{\raggedleft\arraybackslash}p{0.41\textwidth} @{}}
\textbf{\#} & \textbf{Proof} & \textbf{Mathematics} \\
\midrule
\textbf{1.} & We prove that proposition 14: the complements about the diameter of a gnomon are equal for Euclid Book II the Euclidean plane. & \\[0.45em]
\textbf{2.} & Let gnomon = gnomon\_about\_diameter. & \\[0.45em]
\textbf{3.} & Let comp\_1 = complement\_about\_diameter\_1. & \\[0.45em]
\textbf{4.} & Let comp\_2 = complement\_about\_diameter\_2. & \\[0.45em]
\textbf{5.} & From step 2, step 3, and step 4, we can deduce that comp 1 equals comp 2. & $\displaystyle \text{comp 1} = \text{comp 2}$ \\[0.45em]
\textbf{6.} & We invoke the derived fact governing Book I of the Elements in the Euclidean plane: proposition 34: complements of parallelograms about the diameter are equal, for Book I of the Elements in on the Euclidean plane. & \\[0.45em]
\textbf{7.} & We invoke the derived fact governing Euclid Book II the Euclidean plane: proposition 11: a gnomon equals the rectangle by the segments, for Euclid Book II on the Euclidean plane. & \\[0.45em]
\textbf{8.} & From step 6 and step 7, this implies complement 1 equals complement 2. Hence proven. & $\displaystyle \text{complement 1} = \text{complement 2}$ \\[0.45em]
\bottomrule
\end{tabular}
\label{thm:Geometry-Euclid Book II-proposition_14_the_complements_about_the_diameter_of_a_gnomon_are_equal}
\par\medskip\hrule
\end{document}
